% BHU PYQ -- Manually transcribed & error-corrected
% Paper: BCH-215 : Fundamentals of Business Finance
% Exam: B.Com. (Hons.) Semester III Examination, 2016-17

\documentclass[11pt,a4paper]{article}
\usepackage[a4paper,top=2.5cm,bottom=2.5cm,left=2.8cm,right=2.8cm,headheight=14pt]{geometry}
\usepackage{amsmath,amssymb}
\usepackage[T1]{fontenc}
\usepackage[utf8]{inputenc}
\usepackage{lmodern,microtype,fancyhdr,enumitem,hyperref,lastpage,parskip}

\hypersetup{
    colorlinks=true,
    urlcolor=black,
    linkcolor=black,
    pdftitle={BCH-215: Fundamentals of Business Finance -- BHU B.Com. (Hons.) Sem III 2016-17}
}

\pagestyle{fancy}
\fancyhf{}
\renewcommand{\headrulewidth}{0.4pt}
\renewcommand{\footrulewidth}{0.4pt}
\fancyhead[L]{\small\textit{Banaras Hindu University}}
\fancyhead[R]{\small\textit{BCH-215: Fundamentals of Business Finance}}
\fancyfoot[C]{\small Page \thepage\ of \pageref{LastPage}}

\newlist{parts}{enumerate}{2}
\setlist[parts,1]{label=(\alph*),leftmargin=*,itemsep=5pt,topsep=3pt}
\setlist[parts,2]{label=(\roman*),leftmargin=*,itemsep=3pt,topsep=2pt}
\newcommand{\pts}[1]{\hfill\textit{[#1]}}

\begin{document}

\begin{center}
  {\large\bfseries BANARAS HINDU UNIVERSITY}\\[2pt]
  {\normalsize B.Com. (Hons.) --- Semester III Examination, 2016-17}\\[6pt]
  \rule{0.8\linewidth}{0.5pt}\\[6pt]
  {\large\bfseries Paper No. BCH-215: Fundamentals of Business Finance}\\[4pt]
  {\normalsize Time: 3 Hours\hfill Maximum Marks: 70}
\end{center}

\rule{\linewidth}{0.4pt}

\smallskip
\noindent\textit{Instructions: Answer all questions. All questions carry equal marks (14 marks each).}

\bigskip

\noindent\textbf{Question 1.}
Define business finance. ``Wealth maximization is superior to profit maximization.'' Discuss.\pts{14}

\centerline{\textbf{OR}}

State the various functions of finance.\pts{14}

\medskip\rule{\linewidth}{0.2pt}

\noindent\textbf{Question 2.}
What do you mean by a financial plan? State its determinants.\pts{14}

\centerline{\textbf{OR}}

What considerations would you keep in mind while drafting a financial plan for an industrial concern? Also discuss the limitations of a financial plan.\pts{14}

\medskip\rule{\linewidth}{0.2pt}

\noindent\textbf{Question 3.}
Define overcapitalisation. Discuss its causes and impact on shareholders, the company, and society.\pts{14}

\centerline{\textbf{OR}}

Write short notes on the following:\pts{14}
\begin{parts}
  \item Cost theory of capitalisation
  \item Earnings theory of capitalisation
\end{parts}

\medskip\rule{\linewidth}{0.2pt}

\noindent\textbf{Question 4.}
Examine briefly the various sources of long-term capital in Indian industries.\pts{14}

\centerline{\textbf{OR}}

Write short notes on the following:\pts{14}
\begin{parts}
  \item Retained earnings
  \item Equity share capital
\end{parts}

\medskip\rule{\linewidth}{0.2pt}

\noindent\textbf{Question 5.}
What do you mean by working capital? State the factors determining the working capital requirements.\pts{14}

\centerline{\textbf{OR}}

Write short notes on the following:\pts{14}
\begin{parts}
  \item Commercial papers
  \item Trade financing
\end{parts}

\vfill
\rule{\linewidth}{0.4pt}
\begin{center}\small
\href{https://bhu-exam.vercel.app/}{Click here to visit our website}\\[4pt]
\href{https://me-aryan.vercel.app/}{Click here to visit our contributor}\\[4pt]
\href{https://quashan-me.vercel.app/}{Click here to visit our contributor}
\end{center}

\end{document}
